% Pursuing Stacks (cotes 134-2 to 134-8): archive cote 134-8, site page 19 (PDF page 20), typescript page 573.
% typescript. Draft compared: ps7.tex:876-933.
% AI-transcribed by Claude Opus 5.5 (claude-opus-5-5, adaptive thinking, effort high), Message Batches API, batch ps4 (msgbatch_012Aoc1NuskATPsyBeKSYK7J), 2026-11-12; not yet checked by eye.
\scanpage{134-8/19}
\mspage{573}

these,  besides the relevant base change properties,  are so precise,
one feels,  that they may almost be taken as a definition!  Maybe even
the ``almost'' could be dropped  -  namely that a comprehensive axiomatic
set up for the duality formalism could be worked out,  in a way applicab
le to the known instances as well as to the present\xout{\unsure{a}}ly still unknown
one of (Cat), \inserted{by} going a little further still than Deligne's exposition
in SGA 5 (where the notions of ``smooth'' and ``proper'' maps were supposed
to be given beforehand,  satisfying suitable properties). Before diving
into such axiomatization game,  one should get a better feeling though,
through a fai\unsure{r|t} number of examples (not all with $e$ as the target category
moreover),  of how the proper,  the smooth and the proper-and-smooth
maps in (Cat) actually look like.  Here presumably Verdier's work in
the context of discrete infinite groups should give useful clues.

Other important clues should come from the opposite side so to say
- namely ordered sets.  Such \inserted{a} set\xout{s} \inserted{I},  besides defining in the usual way
a \inserted{small} category and hence a topos, \inserted{may} equally \xout{define} be viewed as a topological
space,  more accurately the topos it defines may be viewed as being asso
ciated to a topological space,  admitting I as it's underlying set
(cf. section 22, p.18).  This topological space is noetherian iff the
ordered set I is -  for instance if I is finite.  In such a case, an \xout{al}
old algebraic geometer like me will feel in known territory,  which
maybe is a delusion however - at any rate,  I doubt the duality formali-%
sm for topological spaces (using factorizations of maps $X\hand{\longrightarrow}Y$ via embed-%
dings in \over{S}{\unsure{s}}paces $Y\times\underline{R}^{d}$ ) makes much sense for such non-separated spaces.
However,  as we saw in section 22,  when I satisfies some mild \inserted[typed]{\hand{``}local} finite-%
ness'' requirement (for instance when I is finite),  we may associate to
it a \uline{geometrical realization} \hand{$|$}I\hand{$|$} ,  which is a locally compact space
(a compact one indeed if I is finite) endowed with a \hand{``}conical \xout{cellular}
subdivision\hand{''} \hand{(}indexed by \inserted{the opposite ordered set} I$^{\hand{o}}$\hand{)} hence canonically triangulated via the \hand{``}barycen-%
tric subdivision\hand{''}\deleted{)}.  The homotopy type of this space is canonically iso-%
morphic to the homotopy type of I,  viewed as a ``model'' in (Cat).  What
is more important here,  is that a \inserted{(pre)}sheaf of sets (say) on the category
I  may be interpreted as being essentially the same as a sheaf of sets
on the geometric realization \hand{$|$}I\hand{$|$} \uline{which is locally constant} (\uline{and} \uline{hence}
\uline{constant}) \uline{on each of the} ``open'' \uline{strata or \hand{``}cones\hand{''} of}  I .  This descripti-%
on then carries over to sheaves of $k$-modules,  The ``clue'' I had in mind
is that \uline{within the context of} \xout{locally} \inserted{locally} \uline{finite ordered sets}\hand{,} \deleted{say,}  the
\uline{looked for} \uline{``six operations duality formalism''} \uline{should be no more,} \uline{no less}
\uline{than the accurate reflection of the same formalism within the context}

\margin{19}
